% !PRINTED 714
\weberpage{714}
But if $\lambda=1$, then $\beta=0$ must hold, and two further
cases are still possible:
\[
\begin{aligned}
&\text{either}\quad f=1,\ \alpha=0,\ p\equiv\phantom{-}1\pmod{m},\ \lambda=f\\
&\text{or}\qquad\quad\, f=2,\ \alpha=1,\ p\equiv -1\pmod{m},\ \lambda=\tfrac12 f.
\end{aligned}
\]
Only in this last case does $p$ belong to the first kind, and it
follows therefore that here too $\Theta^{\beta}$ is in general a
primitive $f_1$th root of unity. With this the theorem is then
proved.

This theorem permits us the application of the formula \S.\,193, (9),
which here gives:
\weq{8}{(1-p^{-sf_1})^{e_1} = \prod^{\chi_1}\bigl[1-\chi_1(p)p^{-s}\bigr],}
in which the product extends over all characters $\chi_1$ of the
group $C_1$, and now the function $\Phi_1(s)$ can be further
developed just as $\Phi(s)$ was in \S.\,193. One obtains, in the
same way as there, the formula (11):
\weq{9}{\Phi_1(s) = \frac{1}{1-q^{-s}}\;\prod^{\chi_1}\sum^{n}\frac{\chi_1(n)}{n^s},}
and when one passes to the limit $s=1$,
\weq{10}{g_1 h_1 = \prod^{\chi_1}\sum^{n}\frac{\chi_1(n)}{n},}
in which, however, the product is now to be extended over all
characters of the group $C_1$ \gesperrt{with the exception of the principal
character}. This product is therefore a part of the product by which
the number $gh$ is expressed in \S.\,193.

The numbers $g$, $g_1$ follow from the general expression given
in \S.\,190, (11):
\weq{11}{\frac{2^{\nu}\pi^{n-\nu}L}{w\sqrt{\pm\Delta}}.}
Since, as was shown in \S.\,178, the units in $\Omega_m$ and
$H_m$, apart from the roots of unity, are the same, a fundamental
system of units in $H_m$ is at the same time a fundamental system
in $\Omega_m$.

If therefore $\varepsilon_1,\varepsilon_2,\dots,\varepsilon_{\nu-1}$
denotes a fundamental system of real units, and if
$\varepsilon_{i,1},\varepsilon_{i,2},\dots,\varepsilon_{i,\nu}$ are
the conjugate units to $\varepsilon_i$, then the conjugate
logarithms in the field $H_m$ are:
\[ \log|\varepsilon_{i,1}|,\ \log|\varepsilon_{i,2}|,\ \dots,\ \log|\varepsilon_{i,\nu}|, \]
whereas in the field $\Omega_m$, in which only conjugate imaginary
pairs occur (\S.\,185),
\[ 2\log|\varepsilon_{i,1}|,\ 2\log|\varepsilon_{i,2}|,\ \dots,\ 2\log|\varepsilon_{i,\nu}|. \]
